% geometry/curvature_classes.tex

\chapter{Curvature Classes}

\label{chap:curvature-classes}

% ============================================================
\section*{Motivation}
% ============================================================

Not all curvature laws possess the same geometric, dynamic, numerical,
or physical properties.

Within AIM, curvature laws are therefore classified according to:
\begin{itemize}
\item continuity,
\item differentiability,
\item boundedness,
\item transition behaviour,
\item realization behaviour,
\item runtime stability,
\item and operational suitability.
\end{itemize}

This classification forms the basis for:
\begin{itemize}
\item transition evaluation,
\item runtime robustness,
\item realization analysis,
\item admissibility checks,
\item and optimization strategies.
\end{itemize}

\begin{principle}[Curvature classification principle]
\label{princ:curvature-classification}
Within AIM, railway alignment elements are classified primarily by their
curvature behaviour rather than by their visual position-space shape.
\end{principle}

% ============================================================
\section{Basic Curvature Classes}
% ============================================================

\subsection{Constant Curvature}

\begin{definition}[Constant curvature]
\label{def:constant-curvature-class}
A curvature law is constant if:
\[
\kappa(s)=\kappa_0.
\]
\end{definition}

Special cases are:
\begin{itemize}
\item \(\kappa_0=0\): straight line,
\item \(\kappa_0\neq 0\): circular arc.
\end{itemize}

Constant curvature forms the fixed-element basis of classical railway
geometry.

% ------------------------------------------------------------

\subsection{Piecewise Constant Curvature}

\begin{definition}[Piecewise constant curvature]
\label{def:piecewise-constant-curvature}
A curvature law is piecewise constant if:
\[
\kappa(s)
=
\kappa_i
\]
on subintervals \(I_i\).
\end{definition}

This corresponds to classical tangent--arc railway geometry.

Piecewise constant curvature is geometrically simple, but curvature
jumps are dynamically and operationally problematic.

% ============================================================
\section{Continuity Classes}
% ============================================================

\subsection{\texorpdfstring{$C^0$}{C0} Curvature}

\begin{definition}[\(C^0\)-continuous curvature]
\label{def:c0-curvature}
A curvature law belongs to \(C^0\) if:
\[
\kappa(s)
\]
is continuous.
\end{definition}

Continuous curvature eliminates curvature jumps.

% ------------------------------------------------------------

\subsection{\texorpdfstring{$C^1$}{C1} Curvature}

\begin{definition}[\(C^1\)-continuous curvature]
\label{def:c1-curvature}
A curvature law belongs to \(C^1\) if:
\[
\kappa'(s)
\]
exists and is continuous.
\end{definition}

This improves dynamic smoothness, realization behaviour, and numerical
robustness.

% ------------------------------------------------------------

\subsection{\texorpdfstring{$C^2$}{C2} Curvature}

\begin{definition}[\(C^2\)-continuous curvature]
\label{def:c2-curvature}
A curvature law belongs to \(C^2\) if:
\[
\kappa''(s)
\]
exists and is continuous.
\end{definition}

Higher continuity reduces excitation effects, numerical instability, and
unwanted high-frequency behaviour.

In AIM, continuity class is not merely a mathematical property. It is an
indicator for dynamic quality, runtime stability, and realization
robustness.

% ============================================================
\section{Boundedness Classes}
% ============================================================

\subsection{Bounded Curvature}

\begin{definition}[Bounded curvature]
\label{def:bounded-curvature}
A curvature law is bounded if:
\[
|\kappa(s)|
\le
\kappa_{\max}.
\]
\end{definition}

Bounded curvature corresponds to a minimum admissible radius.

% ------------------------------------------------------------

\subsection{Bounded Curvature Gradient}

\begin{definition}[Bounded curvature gradient]
\label{def:bounded-curvature-gradient}
A curvature law has bounded curvature gradient if:
\[
|\kappa'(s)|
\le
\eta_{\max}.
\]
\end{definition}

This is closely related to jerk limitations in railway dynamics.

% ------------------------------------------------------------

\subsection{Bounded Higher Derivatives}

\begin{definition}[Bounded higher derivatives]
\label{def:bounded-higher-curvature-derivatives}
Higher-order boundedness conditions include:
\[
|\kappa''(s)|<\infty,
\qquad
|\kappa'''(s)|<\infty.
\]
\end{definition}

Higher-order boundedness becomes relevant for realization behaviour,
vehicle response, and transition-family comparison.

% ============================================================
\section{Transition Classes}
% ============================================================

\subsection{Monotone Transitions}

\begin{definition}[Monotone transition]
\label{def:monotone-transition}
A transition curvature law is monotone if:
\[
\kappa'(s)
\]
does not change sign.
\end{definition}

Monotone transitions represent direct curvature change between two
curvature states.

% ------------------------------------------------------------

\subsection{Symmetric Transitions}

\begin{definition}[Symmetric transition]
\label{def:symmetric-transition}
A normalized transition law is symmetric if its curvature evolution is
symmetric with respect to the midpoint of the transition interval.
\end{definition}

For normalized transition coordinates \(w\in[0,1]\), this may be
expressed by a symmetry relation of the transition shape function.

Symmetry is not an absolute requirement, but it provides a useful
classification property.

% ------------------------------------------------------------

\subsection{Asymmetric Transitions}

\begin{definition}[Asymmetric transition]
\label{def:asymmetric-transition}
A transition law is asymmetric if symmetry is intentionally violated.
\end{definition}

Asymmetric transitions may appear in:
\begin{itemize}
\item constrained engineering situations,
\item realization-aware optimization,
\item vehicle-dynamics optimization,
\item and topology-constrained alignments.
\end{itemize}

% ============================================================
\section{Analytical Classes}
% ============================================================

\subsection{Polynomial Curvature Laws}

\begin{definition}[Polynomial curvature]
\label{def:polynomial-curvature}
A polynomial curvature law satisfies:
\[
\kappa(s)
=
\sum_{i=0}^n a_i s^i.
\]
\end{definition}

Polynomial curvature laws are useful because derivatives and integrals
can often be handled systematically.

% ------------------------------------------------------------

\subsection{Trigonometric Curvature Laws}

\begin{definition}[Trigonometric curvature]
\label{def:trigonometric-curvature}
A trigonometric curvature law contains sinusoidal or cosine-based
components.
\end{definition}

Trigonometric curvature laws are useful for smooth transition behaviour
and frequency-oriented interpretation.

% ------------------------------------------------------------

\subsection{Lookup-Based Curvature Laws}

\begin{definition}[Lookup-based curvature]
\label{def:lookup-based-curvature}
A lookup-based curvature law is defined by sampled or tabulated
representations together with interpolation rules.
\end{definition}

Lookup-based laws allow practical inclusion of transition families or
realization effects that are not conveniently represented by a simple
closed formula.

% ------------------------------------------------------------

\subsection{Hybrid Curvature Laws}

\begin{definition}[Hybrid curvature law]
\label{def:hybrid-curvature-law}
A hybrid curvature law combines:
\begin{itemize}
\item parametric structure,
\item sampled corrections,
\item measured data,
\item and local refinement.
\end{itemize}
\end{definition}

Hybrid laws are especially important for measured and realized railway
states.

% ============================================================
\section{Sparse Curvature Classes}
% ============================================================

\begin{definition}[Sparse curvature representation]
\label{def:sparse-curvature-class}
A sparse curvature model represents geometry through:
\begin{itemize}
\item fixed curvature elements,
\item transition elements,
\item and compact parameter sets.
\end{itemize}
\end{definition}

Sparse classes emphasize:
\begin{itemize}
\item compactness,
\item runtime efficiency,
\item reconstruction stability,
\item and semantic clarity.
\end{itemize}

Within AIM, sparse curvature classes are not merely storage formats.
They define executable curvature programs.

% ============================================================
\section{Operational Classes}
% ============================================================

\subsection{Dynamically Smooth Curvature}

\begin{definition}[Dynamically smooth curvature]
\label{def:dynamically-smooth-curvature}
A curvature law is dynamically smooth if:
\begin{itemize}
\item jerk remains bounded,
\item cant-rate remains admissible,
\item curvature variation is controlled,
\item and high-frequency excitation is limited.
\end{itemize}
\end{definition}

% ------------------------------------------------------------

\subsection{Realization-Stable Curvature}

\begin{definition}[Realization-stable curvature]
\label{def:realization-stable-curvature}
A curvature law is realization-stable if:
\begin{itemize}
\item realization preserves its dominant structure,
\item construction processes do not introduce instability,
\item and its low-frequency behaviour remains identifiable after
physical realization.
\end{itemize}
\end{definition}

Realization-stable curvature is a physical property, not merely an
analytical smoothness property.

% ------------------------------------------------------------

\subsection{Runtime-Stable Curvature}

\begin{definition}[Runtime-stable curvature]
\label{def:runtime-stable-curvature}
A curvature law is runtime-stable if:
\begin{itemize}
\item projection remains numerically stable,
\item frame evaluation remains continuous,
\item pose reconstruction behaves robustly,
\item and local inverse problems remain well-conditioned.
\end{itemize}
\end{definition}

% ============================================================
\section{Frequency Classes}
% ============================================================

\subsection{Low-Frequency Curvature}

\begin{definition}[Low-frequency curvature]
\label{def:low-frequency-curvature}
Low-frequency curvature components determine large-scale geometric
behaviour.
\end{definition}

These components are usually preserved by construction and maintenance
processes.

% ------------------------------------------------------------

\subsection{High-Frequency Curvature}

\begin{definition}[High-frequency curvature]
\label{def:high-frequency-curvature}
High-frequency curvature components influence:
\begin{itemize}
\item realization sensitivity,
\item dynamic roughness,
\item measurement sensitivity,
\item and numerical instability.
\end{itemize}
\end{definition}

\begin{proposition}[Spectral realization principle]
\label{prop:curvature-classes-spectral-realization}
High-frequency curvature components are attenuated by realization.
\end{proposition}

This frequency-oriented classification connects curvature classes with
the realization filter interpretation.

% ============================================================
\section{Admissibility}
% ============================================================

\begin{definition}[Admissible curvature law]
\label{def:admissible-curvature-law}
A curvature law is admissible if it satisfies:
\begin{itemize}
\item geometric constraints,
\item realization constraints,
\item operational constraints,
\item dynamic constraints,
\item and numerical stability conditions.
\end{itemize}
\end{definition}

Admissibility depends on:
\begin{itemize}
\item vehicle dynamics,
\item realization processes,
\item standards,
\item topology,
\item runtime requirements,
\item and intended operational use.
\end{itemize}

Thus, admissibility is context-dependent.
A curvature law may be analytically valid but operationally unsuitable
or physically unstable under realization.

% ============================================================
\section{Classification versus Quality}
% ============================================================

Classification does not automatically imply quality.

A highly smooth curvature law may still be unsuitable if it is:
\begin{itemize}
\item difficult to realize,
\item dynamically unfavourable,
\item numerically unstable,
\item incompatible with topology,
\item or unsuitable for the intended operation.
\end{itemize}

Within AIM, curvature quality is evaluated by the interaction of class,
context, realization, and runtime behaviour.

% ============================================================
\section{Canonical Interpretation}
% ============================================================

\begin{proposition}
Different curvature classes correspond to fundamentally different
geometric, operational, runtime, and realization behaviours.
\end{proposition}

\begin{proposition}
Within AIM, railway alignment quality is primarily characterized in
curvature space rather than position space.
\end{proposition}

\begin{proposition}
Transition families are interpreted as structured subsets of curvature
classes.
\end{proposition}

\begin{proposition}
Admissibility is not a purely geometric property, but a combined
geometric, dynamic, runtime, and realization property.
\end{proposition}

% ============================================================
\section{Connection to AIM}
% ============================================================

\begin{summary}
Curvature classes provide the structural classification framework of AIM
geometry.

They connect:
\begin{itemize}
\item continuity,
\item differentiability,
\item boundedness,
\item transition behaviour,
\item realization behaviour,
\item runtime stability,
\item dynamic smoothness,
\item and optimization suitability.
\end{itemize}

Within AIM, transition curves are therefore interpreted not merely as
geometric entities, but as members of specific curvature classes with
distinct operational properties.

This creates the bridge from intrinsic geometry to classification,
realization analysis, runtime evaluation, and optimization.
\end{summary}
